The cold had found his hand first, as it always did now, and it had taken Jaime the better part of a fortnight on the kingsroad to understand why.

A living hand told a man things. It ached before a storm, stiffened in true cold, warned him in a hundred small ways before the weather turned bad enough to matter. The gold hand told him nothing at all, until the nothing itself became a kind of message — a dead weight strapped to the end of his arm, indifferent to frost the way it was indifferent to everything, and it was only on the ninth or tenth day north of the causeway, when he noticed the stump above it had gone the deep aching red of a burn, that he understood the metal was drawing the warmth out of him rather than simply failing to provide any. Cold country wanted a man’s whole body to fight it. His arm was fighting a war on behalf of ornamental metalwork he’d once been vain enough to think looked well on him.

He had stopped removing the glove entirely by the time the trees started closing in around the kingsroad, the deep dark pines the Northmen called the wolfswood pressing close enough on either side that the sky narrowed to a grey seam far overhead. He rode with the stump wrapped in wool inside the glove now, which helped some, and thought, not for the first time on this ride, that there was something fitting in it — that the hand Aerys Targaryen’s death had cost him should turn out to be the one part of him least suited to surviving the second war he’d ridden north to end.

He had not told anyone he was leaving alone. That had not been the plan, once, back when the plan had still involved marching north at the head of the army he’d spent a year building for her — for Cersei, who had promised Winterfell an army and never intended to send one, who had smiled at him across the map table with the particular warmth she reserved for lies she expected him to believe, and had believed, gods help him, right up until the moment he understood he was the only person in that room who still did.

\textit{You’re leaving,} she’d said, when he told her, and it had not been a question, and there had been no fear in it at all, only a kind of contempt so total it had almost circled back around into something like respect — as if some small part of her had always known this was coming and had simply been waiting to see how long it would take him.

\textit{Someone from our house should be there,} he’d told her. \textit{When it comes. Someone should stand for something.}

\textit{Stand for what? For a dead man’s army walking on Sansa Stark’s snow? Go and stand for it, then, and see how warmly they welcome the Kingslayer to their war.}

She had not asked him to stay. That was the thing he kept returning to, riding, the wolfswood pressing in dark and silent on either side of him — that in all their years, all their children, all the things they had done and lied about and lived through together, she had never once, in the end, asked him to stay for her own sake. Only for the babe’s. And even that she had said the way she said everything now, as if daring him to notice how little it cost her to say it.

He rode alone because there had been no one left he trusted to ride with, and because some part of him — the part still capable of shame, which he had thought Cersei burned out of him years before and which kept, infuriatingly, reasserting itself at the worst possible times — did not think he deserved company for this particular ride. He was going north to stand in a hall full of people who had every reason to want him dead, to face a boy he had thrown from a tower window and crippled for life, to offer himself to a war he had spent most of his life assuming was a story soldiers told new recruits to keep them afraid of the dark.

He was three days from Winterfell, by his best reckoning, when he stopped believing it was a story.

\sceneBreak

The steading had perhaps housed thirty souls, going by the size of it — a scatter of low turf-roofed houses around a well, the kind of place that would not have a name most maps troubled to record. He smelled it before he saw it, the particular sourness of a fire long dead and something else beneath the smoke that he made himself not name until he’d seen it directly, because a man could sometimes talk himself out of a smell if he didn’t look too hard at what made it.

He looked hard. There was no talking himself out of it.

Whatever had come through had not troubled to hide its work. The bodies — those left whole enough to count as bodies — lay where they had fallen, in doorways, across the well’s stone lip, one small shape that had been a child not moved at all from where it must have been playing when whatever came had come. There were no wounds on most of them that Jaime could find, no blood in the quantities a blade or an axe left behind. They looked, more than anything, like people who had simply stopped, mid-motion, the way a puppet stops when the hand inside it is withdrawn.

None of them had risen again to walk. That, at least, was mercy of a kind — whatever host had passed through here had not troubled to make soldiers of thirty dead villagers, either because there was no need of them yet or because the making required something Jaime did not understand and did not want to.

He found the one living soul by her coughing before he found her by sight — an old woman, wrapped in three cloaks that were not enough against a cold that had gotten inside her days ago by the sound of her breathing, hidden in the root cellar beneath what had been, he supposed, the steading’s largest house. She came up fighting when he pulled the cellar door back, an old kitchen knife in a hand shaking too badly to do murder with it, and stopped only when she saw the white cloak, filthy as it was, and the golden hand, and worked out, faster than he’d have credited, exactly who was kneeling in front of her.

“Kingslayer,” she said. It was not a question either. Northerners, he was coming to understand, did not ask him things. They simply named him and waited to see what he’d do about it.

“Aye.” There seemed no profit in denying it, this far past denying anything. “How long since they came through?”

“Four days. Five.” Her voice was a dry husk of itself. “They don’t sleep, ser. They don’t stop. My boy went out to bar the door and I never — “ She did not finish it, and Jaime did not make her.

He gave her what food he carried, and his second cloak, and walked her the half mile to where the kingsroad rejoined itself past the steading’s boundary stones, because there was nothing else in that place worth staying for and a woman that cold would not survive a night alone in it. She did not thank him. He had not expected her to. Somewhere around the third mile, walking his horse rather than riding it so she could lean on the beast’s flank for warmth, she said, without looking at him, “My boy would’ve spat on you a fortnight past. For what you did to the old king.” A pause, her breath fogging thin in the cold. “Don’t seem to matter much now, does it. What anyone did to anyone. Not with that coming.”

“No,” Jaime agreed. “I don’t suppose it does.”

It was, he thought, riding on alone once he’d left her with a crofter’s family two miles further who took her in without much discussion at all, the closest thing to absolution anyone in this country was likely to offer him, and he understood, turning it over through the following short grey day, that it would have to be enough, because it was all there was.

\sceneBreak

Winterfell’s walls came up grey and vast out of a grey sky, snow already beginning to fall in the thin, serious way that promised more behind it, and Jaime sat his tired horse at the edge of the tree line for a long moment before he made himself ride on, because he had discovered, somewhat to his own surprise, that the last hundred yards of a very long ride could be the hardest part of it.

The gate guards saw the white cloak first — filthy, travel-stained, but unmistakable even at a distance, the one garment in Westeros that announced a man’s history before he’d said a word of it. Two spears came down to bar the road before he’d covered half the distance, and a third man, older, a scarred Northman with a raven badge that marked him castellan of the gate rather than common guard, stepped out in front of the both of them with the flat, unhurried authority of a man who had turned away worse than one rider in his time.

“Hold there.” No fear in it. Just the North’s particular species of contempt, plain and unhurried as the snow. “State yourself.”

“Jaime Lannister.” No point in the rest of it — the guard’s eyes had already gone to the hand, to the cloak, and settled into the look Jaime had come to recognize as well as his own reflection: recognition arriving a half-breath before understanding did, and understanding arriving a half-breath before the man decided what to feel about it. “I’ve come to fight.”

“Have you.” The guard did not lower his spear. Behind him, up along the battlements, Jaime could see other faces turning, word already beginning to move along the wall the way word always did, faster than any rider could outpace it. “Kingslayer, come alone, three days ahead of any army, to fight in a war his own sister swore she’d send men for.” The scarred Northman looked him over once more, unhurried, the way a man looks over a horse he suspects has been lamed and sold to him anyway. “We’ll want to know a great deal more than that before you’re inside these walls, ser. Starting with where the army you swore us actually is.”

“There isn’t one,” Jaime said, and felt the words land in his own chest with all the weight he’d been trying to outrun the length of the wolfswood. “There’s only me.”

The snow went on falling, and the guard’s spear did not move, and somewhere behind the gate, faint but carrying, Jaime heard a woman’s voice — cool, controlled, unmistakably a Stark’s voice even before he had a name to put to it — asking what in the seven hells was taking so long at her gate.